% =====================================================================
%  THÈME 2 — Notation scientifique & unités — EXERCICES — Niveau DIFFICILE
% =====================================================================
\documentclass[11pt,a4paper]{article}
\usepackage[utf8]{inputenc}
\usepackage[T1]{fontenc}
\usepackage{lmodern}
\usepackage[french]{babel}
\usepackage{amsmath,amssymb,mathtools}
\usepackage{icomma}
\usepackage{siunitx}
\sisetup{output-decimal-marker={,},per-mode=symbol}
\usepackage[version=4]{mhchem}
\usepackage{xcolor}
\usepackage[most]{tcolorbox}
\usepackage{enumitem}
\usepackage{fancyhdr}
\usepackage[a4paper,margin=2.2cm,headheight=26pt,headsep=8pt]{geometry}

\definecolor{primary}{RGB}{150,98,162}
\definecolor{primarydark}{RGB}{92,52,110}
\newcommand{\cs}{chiffre significatif}
\newcommand{\CS}{chiffres significatifs}

\pagestyle{fancy}\fancyhf{}
\fancyhead[L]{\small\textcolor{primarydark}{\textbf{Fiche 1} — Chiffres significatifs}}
\fancyhead[R]{\small\textcolor{primarydark}{\textsc{Exercices · Difficile · Thème 2}}}
\fancyfoot[C]{\small\thepage}
\renewcommand{\headrulewidth}{0.8pt}
\renewcommand{\headrule}{\hbox to\headwidth{\color{primary}\leaders\hrule height \headrulewidth\hfill}}

\newtcolorbox[auto counter]{exo}[1][]{enhanced,breakable,boxrule=0.8pt,arc=2pt,
  colback=primary!4,colframe=primary,left=3mm,right=3mm,top=2mm,bottom=2mm,
  fonttitle=\bfseries,coltitle=white,
  title={Exercice~\thetcbcounter\ifx\relax#1\relax\relax\else\ ---\ \textmd{\itshape #1}\fi}}

\setlength{\parindent}{0pt}\setlength{\parskip}{3pt}

\begin{document}

\begin{center}
{\Large\bfseries\color{primarydark} Exercices — Niveau Difficile}\\[1pt]
{\color{primary} Thème 2 : notation scientifique \& changements d'unité}
\end{center}
\vspace{1mm}
\textit{Chaînes de conversions, nombres exacts et notation d'ingénieur. Traite les sous-questions
dans l'ordre.}
\vspace{2mm}

\begin{exo}[Chaîne de conversions à précision constante]
Une dose de principe actif vaut $m = \num{4,50e2}\,\si{\milli\gram}$.
\begin{enumerate}[label=\alph*),leftmargin=1.6em,itemsep=1pt]
  \item Combien de \CS\ possède $m$ ? Réécris $m$ en écriture décimale usuelle (en mg).
  \item Exprime $m$ en grammes, puis en kilogrammes.
  \item Donne $m$ en kilogrammes sous forme scientifique.
  \item Vérifie que le nombre de CS est le même à chaque étape ; nomme le principe en jeu.
\end{enumerate}
\end{exo}

\begin{exo}[Un nombre exact dans une préparation]
Au laboratoire, on remplit \textbf{3} fioles jaugées identiques, chacune contenant $\qty{25,0}{\milli\litre}$
de solution, et on réunit leur contenu.
\begin{enumerate}[label=\alph*),leftmargin=1.6em,itemsep=1pt]
  \item Parmi les nombres « $3$ » et « $\num{25,0}$ », lequel est exact, lequel est mesuré ? Combien
        de CS pour chacun ?
  \item Calcule le volume total réuni.
  \item Combien de \CS\ le volume total doit-il comporter ? Quelle donnée fixe cette limite ?
  \item En une phrase : pourquoi le facteur « $3$ » ne dégrade-t-il pas la précision ?
\end{enumerate}
\end{exo}

\begin{exo}[Très grand, très petit, et notation d'ingénieur]
\begin{enumerate}[label=\alph*),leftmargin=1.6em,itemsep=1pt]
  \item La concentration en ions oxonium d'une solution vaut $[\ce{H3O+}]=\num{0,00000032}\,\si{\mol\per\litre}$.
        Écris-la en notation scientifique.
  \item Un échantillon contient $\num{8500000}$ cellules. Écris ce nombre en notation scientifique
        (précision : $2$~CS).
  \item La \emph{notation d'ingénieur} n'utilise que des exposants multiples de $3$. Réécris
        $V=\num{2,096e-2}\,\si{\litre}$ en notation d'ingénieur, puis donne l'unité à préfixe SI
        correspondante.
\end{enumerate}
\end{exo}

\begin{exo}[Raisonner par conversions : le mercure]
La masse volumique du mercure vaut $\rho = \qty{13,55}{\gram\per\centi\metre\cubed}$ à \qty{20}{\celsius}.
On rappelle que $\qty{1}{\centi\metre\cubed}=\qty{1}{\milli\litre}$ (exact) et que la masse volumique
de l'eau pure vaut $\qty{1,00}{\gram\per\centi\metre\cubed}$.
\begin{enumerate}[label=\alph*),leftmargin=1.6em,itemsep=1pt]
  \item \textbf{Affirmation A} : « une masse de $\qty{13,55}{\kilo\gram}$ de mercure occupe un volume
        d'exactement $\qty{1}{\litre}$ ». Vérifie-la par conversions (passe par les \si{\centi\metre\cubed}).
  \item \textbf{Affirmation B} : « $\qty{1}{\centi\metre\cubed}$ de mercure a la même masse que
        $\qty{13,55}{\centi\litre}$ d'eau pure ». Vrai ou faux ? Justifie par le calcul.
\end{enumerate}
\end{exo}

\end{document}
